Federated learning lets organisations train a shared intrusion detection model
without exchanging traffic records, but averaging works poorly when
participants hold dissimilar data. Multi-tier personalised methods address this
by placing a group-level model between the device and the server. This work
evaluates one such method, PerMFL, on network intrusion detection data, a
domain it has not been applied to.

The published EMNIST-10 result is reproduced to within 0.04 percentage points.
Sixteen defects are documented in the released implementation, including an
argument binding error that silently overrides a documented hyperparameter with
an unrelated loop count, and a fixed random seed that causes repeated runs to
return identical results.

The method's central architectural feature, the team tier, is then tested
directly. Removing it entirely changes personalised macro F1 by 0.0002 on
CICIDS2017 and by 0.0000 on the base paper's own dataset, against a measured
noise floor of 0.0116. Teams matching true attack domains do not outperform
teams assigned at random. The reason is derived from the update rule: a single
parameter governs both the device-to-team penalty and the weight the team model
places on its members, and at the setting that personalises well, membership
influences the team model by 3.3 per cent. The method's own convergence
condition caps that influence below fifty per cent.

Separating the two roles improves both the personalised and the global model at
every heterogeneity level measured across CICIDS2017, NSL-KDD and TON-IoT,
raising personalised macro F1 by 0.027 to 0.326 and global macro F1 by 0.113 to
0.705, and reducing across-seed variance in global accuracy by two to three
orders of magnitude. One configuration regresses and is reported.
